% !TeX root = ../main.tex
\documentclass[../main]{subfiles}
\begin{document}
\chapter{Inteligencia Artificial y Protocolo MCP en CAD}
\img{images/practicas.jpg}{Integración de IA en el diseño}
\section{Marco Teórico}

La irrupción de la Inteligencia Artificial (IA) ha transformado el Diseño Asistido por Computadora (CAD). Herramientas como el Model Context Protocol (MCP) permiten a los Agentes de IA interactuar de manera nativa con el entorno de diseño.

Un \textbf{Agente de IA} es un sistema autónomo capaz de percibir su entorno, tomar decisiones y ejecutar acciones para lograr objetivos específicos. En el contexto del CAD, estos agentes pueden automatizar tareas repetitivas, optimizar diseños estructurales y sugerir mejoras ergonómicas en tiempo real.

El \textbf{MCP (Model Context Protocol)} es un estándar abierto que facilita la integración entre modelos de lenguaje (LLMs) y herramientas locales o remotas. En el software CAD, el MCP actúa como un puente que permite al agente de IA leer el árbol de diseño, modificar bocetos, ejecutar macros o exportar planos de manera segura y controlada por el usuario. Esto significa que el técnico puede pedir mediante lenguaje natural: ``Crea un engranaje de 50 dientes con módulo 2'', y el agente, mediante MCP, ejecutará los comandos internos del software de CAD para generarlo. 

A diferencia de un comando tradicional escrito manualmente, una instrucción dada mediante lenguaje natural permite generar geometría compleja, optimizar estructuras o sugerir diseños sin necesidad de programar cada paso, interactuando el Agente de IA, el software CAD y el MCP en un flujo de trabajo moderno e integrado. Al ser un estándar abierto, el MCP permite que cualquier desarrollador o técnico adapte y conecte diferentes modelos de IA a distintos programas de CAD, garantizando interoperabilidad.

\section{Cuestionario}
\actividad{Responder las siguientes preguntas basándose estrictamente en el texto anterior.}
\begin{enumerate}
  \item ¿Qué significa la sigla CAD y cómo ha influido la Inteligencia Artificial en este campo?
  \item ¿Qué es un Agente de Inteligencia Artificial (IA) en el contexto del diseño técnico?
  \item Mencione al menos dos tareas que un Agente de IA puede automatizar u optimizar en el software de CAD.
  \item ¿Qué significa la sigla MCP y qué función cumple en el diseño asistido?
  \item ¿Cómo actúa el MCP como puente entre los modelos de lenguaje (LLMs) y el software de CAD?
  \item ¿Qué elementos del diseño puede leer o modificar un agente de IA gracias al MCP?
  \item ¿Cuál es la ventaja de que el agente opere mediante el MCP en lugar de un control externo no integrado?
  \item De acuerdo al texto, ¿qué ejemplo se da de una instrucción mediante lenguaje natural y qué parámetros requiere?
  \item ¿Qué diferencia existe entre un comando tradicional escrito por el usuario y una instrucción dada mediante lenguaje natural a un agente de IA?
  \item Explique cómo interactúan el Agente de IA, el software CAD y el MCP en un flujo de trabajo moderno.
  \item ¿Por qué el texto menciona que el MCP es un ``estándar abierto'' y por qué esto es beneficioso para el entorno técnico?
  \item ¿Qué tipo de objetivos busca lograr un Agente de IA autónomo?
  \item Mencione un beneficio principal de utilizar el Model Context Protocol (MCP) en el flujo de diseño y fabricación de piezas.
\end{enumerate}

\end{document}